\clearpage
\phantomsection
\addcontentsline{toc}{chapter}{Anexos}
\chapter*{Anexos}

\section*{Anexo A. Material complementario del experimento}

Los siguientes archivos, mencionados en el capítulo~\ref{cap:diseno}, se entregan junto con este documento en la carpeta \texttt{anexos/} y permiten reproducir el experimento y el análisis:
\begin{itemize}
  \item \texttt{Bench.java}: banco de pruebas con las dos implementaciones.
  \item \texttt{correr.sh}: \textit{script} de ejecución por bloques.
  \item \texttt{ejecutar.txt}: órdenes utilizadas para lanzar la corrida de 12 bloques.
  \item \texttt{analisis.py}: \textit{script} de análisis estadístico y de generación de gráficas.
  \item \texttt{corrida.log}: registro de la corrida, con la verificación de equivalencia funcional, la calibración y el orden aleatorio de las condiciones en cada bloque.
  \item \texttt{resultados/datos.csv}: una fila por medición (48 corridas y 12 mediciones en reposo).
  \item \texttt{resultados/resumen\_estadistico.csv}: medias, desviaciones, diferencias relativas, intervalos de confianza, \textit{p}-valores y $d$ de Cohen por carga y métrica.
\end{itemize}

\section*{Anexo B. Reproducción del experimento y del análisis}

Con Java y los permisos de lectura de los contadores RAPL, la corrida completa se lanza desde la carpeta \texttt{anexos/} con:

\begin{quote}
\texttt{REPS=12 ./correr.sh}
\end{quote}

El análisis estadístico y las figuras se regeneran a partir de los datos medidos con:

\begin{quote}
\texttt{python3 analisis.py resultados/datos.csv}
\end{quote}

El \textit{script} de análisis requiere Python con las bibliotecas \texttt{numpy}, \texttt{pandas}, \texttt{scipy} y \texttt{matplotlib}.
